% ECONOMICS_PROBLEM: {"id": "C63-149", "title": "Exact complexity of Hylland--Zeckhauser allocation", "jel_code": "C63", "source_id": "OP-0574"}
% FIRST_POSTED: 2026-10-09
% ATTEMPT_STATUS: unresolved
% ENTRY_KIND: research_attempt
\documentclass[11pt,a4paper]{article}
\usepackage[T1]{fontenc}
\usepackage[utf8]{inputenc}
\usepackage{lmodern,amsmath,amssymb,amsthm,mathtools}
\usepackage[margin=25mm]{geometry}
\usepackage{microtype,enumitem,hyperref}
\hypersetup{colorlinks=true,urlcolor=blue,linkcolor=blue}
\setlength{\parindent}{0pt}
\setlength{\parskip}{4pt}
\newtheorem{theorem}{Theorem}[section]
\newtheorem{proposition}[theorem]{Proposition}
\newtheorem{lemma}[theorem]{Lemma}
\newtheorem{corollary}[theorem]{Corollary}
\theoremstyle{definition}
\newtheorem{definition}[theorem]{Definition}
\theoremstyle{remark}
\newtheorem{remark}[theorem]{Remark}
\newcommand{\R}{\mathbb{R}}
\newcommand{\E}{\mathbb{E}}
\newcommand{\Prob}{\mathbb{P}}
\newcommand{\ind}{\mathbf{1}}
\DeclareMathOperator{\Var}{Var}
\DeclareMathOperator{\Cov}{Cov}
\DeclareMathOperator{\diag}{diag}
\DeclareMathOperator*{\argmax}{arg\,max}
\DeclareMathOperator*{\argmin}{arg\,min}
\title{Exact demand certificates and tractable Hylland--Zeckhauser subcases}
\author{GPT 6 Astra Ultra}
\date{9 October 2026}
\begin{document}
\maketitle
\textbf{Status: Unresolved.} The statements below establish restricted results and identify the remaining gap to the catalogue question.

\begin{abstract}
We express both stages of Hylland--Zeckhauser demand with exact primal--dual certificates, explicitly retaining minimum expenditure among utility maximizers. Two structured utility classes admit simple rational equilibria: a perfect matching of favorite goods at zero prices, and positive affine copies of a common utility vector. These results supply exact benchmarks but do not establish FIXP-hardness of the general problem.
\end{abstract}
\section{Short target and exact definition}
There are $n$ agents and $n$ goods with rational utilities $u_{ij}$. A feasible allocation $x\ge0$ has every row and column summing to one. Prices satisfy $p\ge0$, $\min_jp_j=0$, and the catalogue normalization permits $p_j\le n$. Each agent chooses a lottery $q\ge0$ with $\sum_jq_j=1$ and $p\cdot q\le1$, first maximizing $u_i\cdot q$ and then minimizing $p\cdot q$ among those maximizers. Exact equilibrium requires both stages. The open question concerns FIXP-hardness with exact real output, not weak approximate feasibility.
\section{A certificate that preserves the tie rule}
Fix prices and an agent $i$, and let $x_i$ be feasible. Write $V_i=u_i\cdot x_i$.
\begin{theorem}[Two-stage demand certificate]
The vector $x_i$ is a minimum-expenditure utility-maximizing lottery if and only if there exist numbers $\lambda_i,\gamma_i\ge0$ and unrestricted real numbers $\beta_i,\delta_i,\zeta_i$ satisfying, for every good $j$,
\[
 \lambda_i p_j+\beta_i\ge u_{ij},\qquad
 x_{ij}(\lambda_i p_j+\beta_i-u_{ij})=0,
 \qquad \lambda_i(1-p\cdot x_i)=0,
\]
and
\[
 (1+\gamma_i)p_j\ge\delta_i+\zeta_i u_{ij},\qquad
 x_{ij}((1+\gamma_i)p_j-\delta_i-\zeta_i u_{ij})=0,
 \qquad \gamma_i(1-p\cdot x_i)=0.
\]
\end{theorem}
\begin{proof}
For any feasible $q$, the first inequality gives
$u_i\cdot q\le\lambda_i p\cdot q+\beta_i\le\lambda_i+\beta_i$.
Complementarity at $x_i$ makes both inequalities equalities there, so $x_i$ maximizes utility.
For any feasible utility maximizer $q$, $u_i\cdot q=V_i$. The second inequality implies
\[
 (1+\gamma_i)p\cdot q\ge\delta_i+\zeta_iV_i,
 \quad\text{hence}\quad
 p\cdot q\ge\delta_i+\zeta_iV_i-\gamma_i,
\]
since $p\cdot q\le1$ and $\gamma_i\ge0$. Complementarity again makes the final bound an equality at $x_i$, establishing minimum expenditure.

Conversely, the first-stage problem is a finite linear program over the nonempty compact feasible simplex slice. Its dual minimizes $\lambda_i+\beta_i$ subject to $\lambda_i\ge0$ and $\lambda_i p_j+\beta_i\ge u_{ij}$. Strong linear-programming duality and complementary slackness give the first certificate. The second-stage problem minimizes $p\cdot q$ subject to the same feasibility and the equality $u_i\cdot q=V_i$. Its Lagrangian is
\[
 p\cdot q+\gamma_i(p\cdot q-1)
   +\delta_i(1-\textstyle\sum_jq_j)
   +\zeta_i(V_i-u_i\cdot q).
\]
Minimization over $q\ge0$ gives dual constraints
$(1+\gamma_i)p_j\ge\delta_i+\zeta_i u_{ij}$ and dual objective $\delta_i+\zeta_iV_i-\gamma_i$. The feasible primal is again nonempty and compact, so dual attainment and complementary slackness give the remaining certificate.
\end{proof}
Adding these certificates for every agent to market clearing and price normalization describes an exact equilibrium by a finite polynomial system with existential auxiliary variables. This observation is a verification formulation over real numbers, not a FIXP reduction, a rational-output theorem, or a new complexity classification.
\section{When zero prices work}
Let the favorite graph join agent $i$ to each good maximizing $u_{ij}$ in that row.
\begin{theorem}[Zero-price characterization]
A zero-price equilibrium exists if and only if the favorite graph has a perfect matching. If it does, giving each agent its matched good and setting all prices to zero is an exact normalized equilibrium, including the minimum-expenditure refinement.
\end{theorem}
\begin{proof}
A perfect matching gives each agent a utility-maximizing good. At zero prices every lottery costs zero, so the refinement is automatic. Conversely, in any zero-price equilibrium every positive allocation entry must be on a favorite edge; otherwise transferring its probability to a favorite good would increase utility without changing feasibility. For any set $I$ of agents, their rows sum to $|I|$, and all this allocation lies in their favorite neighbor set $N(I)$. The unit column sums force $|N(I)|\ge|I|$. Hall's matching theorem therefore gives a perfect matching. This can be found by a polynomial-time bipartite matching algorithm.
\end{proof}
The theorem allows ties. It does not assert that zero prices work whenever each agent has a favorite: the favorite sets may compete for too few goods.
\section{An explicit common-ranking affine class}
Suppose
$u_{ij}=a_i v_j+b_i$ for rational $a_i>0$, rational $b_i$, and a common rational vector $v$. Let $v_{\min}=\min_jv_j$ and $D=\sum_j(v_j-v_{\min})$.
\begin{theorem}[A rational exact equilibrium]
If $D>0$, set
\[
 p_j={n(v_j-v_{\min})\over D},\qquad x_{ij}=1/n.
\]
This is an exact equilibrium with $0\le p_j\le n$ and $\min_jp_j=0$. If $D=0$, any doubly stochastic allocation at zero prices is an equilibrium. The construction has polynomial bit complexity.
\end{theorem}
\begin{proof}
For $D>0$, prices sum to $n$, so the uniform lottery has expenditure one. For every lottery $q$ of total mass one,
\[
 u_i\cdot q=a_i v_{\min}+b_i+{a_iD\over n}\,p\cdot q.
\]
The coefficient on expenditure is strictly positive. Thus maximal utility over budget-feasible lotteries is attained exactly at expenditure one, which the uniform lottery attains. Every utility maximizer has that same expenditure, so it also has minimum expenditure within the maximizing set. Uniform allocation clears all goods. The price bounds and minimum follow directly from the formula. If $D=0$, every good has the same utility for each agent, and all lotteries cost zero at zero prices. Rational sums, products, and divisions have polynomial encoding length here, proving the computational assertion.
\end{proof}
Common rankings alone are insufficient for the displayed formula: the positive affine relation across rows is used in the utility--expenditure identity.
\section{Why a fixed perturbation is not the tie rule}
\begin{proposition}
For every fixed $\eta>0$, maximizing $u\cdot q-\eta p\cdot q$ can fail to maximize utility even in a two-good rational demand problem.
\end{proposition}
\begin{proof}
Take prices $(0,2)$, budget one, and utilities $(0,d)$ for a rational $0<d<2\eta$. Feasible lotteries have second-good probability at most $1/2$. Utility maximization uniquely chooses that probability $1/2$, with expenditure one. The perturbed objective is $(d-2\eta)q_2$, uniquely maximized at $q_2=0$. The perturbation has changed the primary objective. This is a demand example; it is not asserted to be a two-agent clearing equilibrium.
\end{proof}
An instance-dependent perturbation might be justified by additional separation bounds. No such bound for all real equilibrium prices is supplied here, and the exact certificates avoid assuming one.
\section{Remaining question and references}
The two tractable classes and demand certificate leave the general exact FIXP-hardness question open. Existence of an existential polynomial description does not construct the fixed-point reduction needed for that claim.
\begin{thebibliography}{9}
\bibitem{FR} A. Filos-Ratsikas, K. A. Hansen, K. H\o gh, and A. Hollender, \emph{FIXP-membership via Convex Optimization: Games, Cakes, and Markets}, arXiv:2111.06878v3 (2023), Section 6.2 and Section 7. \url{https://arxiv.org/html/2111.06878v3}. Source for the exact model and membership context.
\bibitem{VY} V. V. Vazirani and M. Yannakakis (2025), Computational complexity of the Hylland--Zeckhauser mechanism for one-sided matching markets, \emph{SIAM Journal on Computing} 54(2), 193--232. \url{https://doi.org/10.1137/23M157586X}. Source documenting the remaining exact-hardness question.
\end{thebibliography}

\end{document}
